% ===== Block B10: 原文 p.82–90 | §4.3 局部自由性 II（一般情形）与因子化 | 译手 tr-ch4 =====
\subsection{局部自由性 II（一般情形）与因子化}
% [原文 p.82]

本节中我们完成层 $\cV_\lambda(\mathfrak F)$ 与 $\cV^{\dagger}_\lambda(\mathfrak F)$ 的局部自由性的证明。

% [Tutor Note: 本节的证明策略先理清：p.82 上半页的推论 4.10 与定理 4.11 已把
%  “纤维光滑的部分” $B\setminus D$ 上的局部自由性解决；本节要处理 $D$ 上（奇异纤维）
%  的情形，办法是每取一个不可约分支 $E=D_i$，把纤维的双点“拆开”（同时正规化）
%  成 $N+2$ 点光滑曲线，再沿 $\sigma',\sigma''$ 用形式参数 $q$ 缝回去。
%  核对：$q\to0$ 时缝合退化为粘点，正对应回到 $E$ 上的奇异纤维。]

\noindent\textbf{4.3.1 缝合.}\quad 设 $\mathfrak F^{(0)}=(\pi:\cC\to\cB;s_1,\dots,s_N)$ 是 $N$-点稳定曲线的一个通用族，而 $\mathfrak F=(\pi:\cC\to\cB;s_1,\dots,s_N;\eta_1,\dots,\eta_N)$ 是 $\mathfrak F^{(0)}$ 上带形式邻域的族。设 $D$ 为族 $\mathfrak F^{(0)}$ 的判别轨迹（discriminant locus），并设 $D=\bigcup_{i=1}^{k}D_i$ 是它分解为不可约分支的分解；这里不可约分支的个数 $k$ 正是该族的纤维的双点个数的最大值。

取定 $D_i$ 并令 $E=D_i$。以 $\pi_E:\cC_E\to E$ 记 $\cC$ 到 $E$ 的限制。设 $\tilde\pi_E:\tilde\cC_E\to E$ 由 $\pi_E:\cC_E\to E$ 沿 $E$ 上方的双点轨迹作同时正规化（simultaneous normalization）得到，而 $\sigma',\sigma'':E\to\tilde\cC_E$ 是相应于 $E$ 上方正规化后的双点的截面。

% [原文 p.83]

\figomitn{$E\to\cB$ 上的交换图：$\tilde\cC_E\to\cC_E\to\cC$，其中 $\pi_E:\cC_E\to E$、$\tilde\pi_E:\tilde\cC_E\to E$ 为投影，$s_j,\sigma$ 与 $s_j,\sigma',\sigma''$ 为相应的截面。}

\begin{proposition}[4.13]
族
\[
 \tilde{\mathfrak F}^{(0)}_E:=(\tilde\pi_E:\tilde\cC_E\to E;\ s_1,\dots,s_N,\sigma',\sigma'')
\]
是 $(N+2)$-点（不必连通）稳定曲线的局部通用族。若 $\tilde\cC_E$ 连通，则 $\tilde{\mathfrak F}^{(0)}_E$ 是亏格 $g-1$ 的稳定曲线族；若不连通，则该族由 $E$ 上两个互不相交的族组成，它们分别是 $E$ 上 $(N'+1)$-点与 $(N''+1)$-点的稳定曲线族，亏格分别为 $g'$ 与 $g''$，其中 $N'+N''=N$ 且 $g'+g''=g$。
\end{proposition}

% [Tutor Note: 命题 4.13 的亏格账可用 Euler 示性数验算：连通情形下把亏格 $g-1$
%  曲线上的两点 $\sigma',\sigma''$ 粘成一个双点，$\chi=2-2(g-1)$ 减少 2，得
%  $2-2g$，正是亏格 $g$；不连通情形把 $g'$ 与 $g''$ 两支上的各一点相粘，得
%  $g'+g''=g$。点数账则为 $(N'+1)+(N''+1)-2=N$，也与原文一致。]

下面考虑族中每条奇异曲线都只有一个双点的情形。此时 $E=D_1$。令 $M=3g-4+2N$。设 $(u_1,u_2,\dots,u_M,q)$ 为 $\cB$ 的坐标，其中 $E$ 由方程 $q=0$ 定义。此外，在 $\cC_E$ 的纤维的双点轨迹 $\Sigma_E$ 的邻域内选取 $\cC_\cB$ 的坐标 $(u_1,\dots,u_M,z,w)$，使得 $\cC_E$ 由 $z=w=0$ 定义，且全纯映射 $\pi$ 由
\begin{equation}
 (u_1,\dots,u_M,z,w)\mapsto(u_1,\dots,u_M,zw) \tag{4.24}
\end{equation}
给出。

% [Tutor Note: $M=3g-4+2N$ 可以这样核对：$\cB$ 的坐标是 $(u_1,\dots,u_M,q)$，
%  共 $M+1=3g-3+2N$ 个；而 $N$-点稳定曲线的模空间维数是 $3g-3+N$，多出的 $N$ 个
%  参数正是每个标记点处形式邻域 $\eta_j$ 的参数（族 $\mathfrak F$ 带形式邻域，
%  见 4.3.1 开头的定义）。总空间 $\cC$ 的局部坐标则是 $M+2$ 个。
%  (4.24) 是结点的标准局部模型 $zw=q$：$q=0$ 时纤维为 $z=0\cup w=0$ 两条分支。]

固定元素 $\mu\in P_\ell$。下面的引理在 2.2 中证明（见引理 2.17）。

\begin{lemma}[4.14]
存在双线性配对
\[
 (\ |\ ):\cH_\mu\times\cH_{\mu^\dagger}\to\C,
\]
它在相差一个常数倍的意义下唯一，并满足
\[
 (X(n)u|v)+(u|X(-n)v)=0
\]
对任意 $X\in\mathfrak g$、$n\in\Z$、$|u\rangle\in\cH_\mu$、$|v\rangle\in\cH_{\mu^\dagger}$ 成立；并且当 $d\ne d'$ 时 $(\ |\ )$ 在 $\cH_\mu(d)\times\cH_{\mu^\dagger}(d')$ 上为零。
\end{lemma}

% [Tutor Note: 这个配对是 $\cH_\mu$ 与 $\cH_{\mu^\dagger}$ 之间的 Shapovalov 型配对，
%  其中 $\mu^\dagger=-w(\mu)$、$w$ 为 Weyl 群最长元（见记号表）。不变性条件写成
%  $(X(n)u|v)=-(u|X(-n)v)$，即配对在 $\hat{\mathfrak g}$ 作用下不变；核对符号：
%  取 $n=0$、$X\in\mathfrak h$ 得 "两边的权相加为零"，与 $\mu^\dagger$ 的定义相符。]

现在选取 $\cH_\mu(d)$ 的基 $\{v_1(d),\dots,v_{m_d}(d)\}$，以及 $\cH_{\mu^\dagger}(d)$ 关于上述双线性型 $(\ |\ )$ 的对偶基 $\{v^1(d),\dots,v^{m_d}(d)\}$。

% [Tutor Note: 对偶基按 $(\ |\ )$ 取，即 $(v_i(d)|v^j(d))=\delta_i^j$（或差一个
%  常数因子，取决于配对归一化）。注意 $v_i(d)$ 与 $v^j(d)$ 分别属于对偶的两个
%  模 $\cH_\mu(d)$ 与 $\cH_{\mu^\dagger}(d)$，不是同一空间的基与对偶基，别误读。]

由推论 4.10，层 $\cV^{\dagger}_{\mu,\mu',\tilde\lambda}(\tilde{\mathfrak F}_E)$ 在 $E\setminus\bigcup_{j\ne i}D_j$ 上是局部自由的。设 $\langle\Psi|$ 是 $\cV^{\dagger}_{\mu,\mu',\tilde\lambda}(\tilde{\mathfrak F}_E)$ 的一个截面。我们定义与 $\langle\Psi|$ 相伴的元素 $\langle\tilde\Psi|\in\cH^{\dagger}_\lambda(E)[[q]]$。为此先定义 $\langle\Psi_d|\in\cH^{\dagger}_\lambda(E)$ 为
\begin{equation}
 \begin{gathered}
 \langle\Psi_d|\Phi\rangle=\sum_{i=1}^{m_d}\langle\Psi|v_i(d)\otimes v^i(d)\otimes\Phi\rangle,\\
 |\Phi\rangle\in\cH^{\dagger}_\lambda(E).
 \end{gathered} \tag{4.25}
\end{equation}
今定义 $\langle\tilde\Psi|\in\cH^{\dagger}_\lambda(E)[[q]]$ 为
\begin{equation}
 \langle\tilde\Psi|\Phi\rangle=\sum_{d=0}^{\infty}\langle\Psi_d|\Phi\rangle q^d. \tag{4.26}
\end{equation}

% [Tutor Note: (4.25) 是“把缝合参数 $q$ 的幂次按 $L_0$-次数 $d$ 逐级算出来”：
%  $\langle\Psi|$ 是 $(N+2)$-点共形块上的线性泛函，要把它收缩成原来 $N$-点问题
%  的泛函，就在 $\sigma',\sigma''$ 处插入 $\cH_\mu\otimes\cH_{\mu^\dagger}$ 的一组
%  对偶基。核对维数：$m_d=\dim\cH_\mu(d)$，求和恰好是沿 $\cH_\mu(d)\otimes\cH_{\mu^\dagger}(d)$
%  的收缩（迹），故 (4.25) 与 (4.26) 是良定义的。]

% [原文 p.84]

\begin{definition}[4.15]
上述带有 $\langle\tilde\Psi|$ 的构造 $q^{\Delta_\mu}\langle\tilde\Psi|$ 称为一个缝合（sewing），其中数 $\Delta_\mu$ 由
\[
 \frac{(\mu,\mu+\rho)}{2(g^*+\ell)}
\]
定义。
\end{definition}

这里数 $g^*$ 是李代数 $\mathfrak g$ 的对偶 Coxeter 数，$\rho$ 在 (2.14) 中定义。

% [Tutor Note: 与记号表对照：$\Delta_\lambda=\frac{(\hat\lambda,\lambda+2\rho)}{2\kappa}$
%  是权 $\lambda$ 的共形权重（$\kappa=g^*+\ell$），而这里 $\Delta_\mu=
%  \frac{(\mu,\mu+\rho)}{2(g^*+\ell)}$，分子中是 $\mu+\rho$ 而非 $\mu+2\rho$，
%  二者一般\textbf{不等}。可验证：取 $\mathfrak g=\mathfrak{sl}_2$、$\mu=\ell\omega_1$，
%  归一化 $(\theta,\theta)=2$ 时 $\Delta_\mu=\frac{\ell(\ell+1)}{4(\ell+2)}$，
%  而共形权重为 $\frac{\ell}{4}$。故按原文抄录，不擅自替换；是否为笔误留待 §5.3 判断。]

在 §5.3 中我们将证明 $\langle\tilde\Psi|$ 确实收敛，且 $q^{\Delta_\mu}\langle\tilde\Psi|$ 定义了共形块层的一个多值截面。

这里我们只证明 $\langle\tilde\Psi|$ 满足形式规范条件。为说明这句话的确切含义，先证明下面的引理。

\begin{lemma}[4.16]
存在 $\cO_E$-模的单射
\[
 \pi_*\cO_{\cC}(*S)|_E\hookrightarrow\tilde\pi_{E*}\cO_{\tilde\cC_E}(*(\sigma'+\sigma''+S))[[q]],
\]
\[
 f\mapsto\sum_{k=0}^{\infty}f_kq^k,
\]
其中
\[
 f_k\in\tilde\pi_{E*}\cO_{\tilde\cC_E}(*S+k(\sigma'+\sigma'')).
\]
\end{lemma}

\noindent\textit{证明.}\quad 取一点 $P\in\cC_E$，它是 $\pi_E$ 的某条纤维的双点。于是如同上面那样，选取以 $P$ 为中心、$\cC$ 的局部坐标 $(u_1,\dots,u_{M-1},z,w)$，以及以 $\pi(P)$ 为中心、$\cB$ 的局部坐标 $(u_1,\dots,u_{M-1},q)$，使得 $\pi$ 由 (4.24) 给出。由于 $f$ 在 $P$ 处全纯，我们有展开
\[
 f=f(u_1,\dots,u_{M_1},z,w)=\sum_{m\ge 0,\ n\ge 0}f_{m,n}(u)z^mw^n.
\]
定义 $g_{P'}(u,q,z)$ 为
\[
 g_{P'}(u,q,z)=f\Bigl(u,z,\frac{q}{z}\Bigr)=\sum_{k=0}^{\infty}g_k(u,z)q^k,
\]
其中
\begin{equation}
 g_k(u,z)=\sum_{m=0}^{\infty}f_{m,k}(u)z^{m-k}. \tag{4.27}
\end{equation}
再定义 $h_{P''}(u,q,w)$ 为
\[
 h_{P''}(u,q,w)=f\Bigl(u,\frac{q}{w},w\Bigr)=\sum_{k=0}^{\infty}h_k(u,w)q^k,
\]
其中
\begin{equation}
 h_k(u,w)=\sum_{n=0}^{\infty}f_{k,n}(u)w^{n-k}. \tag{4.28}
\end{equation}

% [Tutor Note: (4.27) 的推导只有一步代换，务必自己算一遍：把 $w=q/z$ 代入
%  $\sum f_{m,n}z^mw^n$ 得 $\sum_{m,n}f_{m,n}z^{m-n}q^n$，再令 $k=n$，即得
%  $g_k=\sum_m f_{m,k}z^{m-k}$；(4.28) 同理令 $k=m$。这里用到了 $P'$ 与 $P''$
%  分别是双点 $P$ 正规化后位于 $z=0$ 与 $w=0$ 的点，故 $1/z$（或 $1/w$）在对应
%  一侧是极点：$g_k$ 在 $z=0$ 的最高极点为 $z^{-k}$，正与引理中 $f_k$ 允许
%  $k(\sigma'+\sigma'')$ 处的极点吻合。]

% [Tutor Note: 原文 p.84 证明中坐标列表疑有漏字：$\cC$ 处写成
%  $(u_1,\dots,u_{M-1},z,w)$（$M+1$ 个），而 $\dim\cC=\dim\cB+1=M+2$，按 (4.24)
%  前的约定应为 $(u_1,\dots,u_{M-1},q,z,w)$ 或 $(u_1,\dots,u_M,z,w)$；$\cB$ 处
%  $(u_1,\dots,u_{M-1},q)$ 亦比 $\dim\cB=M+1$ 少一个。照原文排，见 ERRATA。]
%  另：同一证明中的展开式显示行原文印作 $f=f(u_1,\dots,u_{M_1},z,w)$，
%  同段行文作 $u_{M-1}$，当为后者（见 ERRATA）。

{\emergencystretch=28.3pt 对于不是纤维的双点的点 $Q\in\cC_E$，我们可以选取以 $Q$ 为中心、$\cC$ 的局部坐标 $(u_1,\dots,u_{M-1},q,z)$，使得 $\pi$ 由到前 $M$ 个因子的投影给出。此时我们有展开\par}
\[
 f(u_1,\dots,u_{M-1},q,z)=\sum_{k=0}^{\infty}f_{Q,k}(u,z)q^k.
\]

% [原文 p.85]

由此可见，数据 $\{g_k(u,z),h_k(u,w),f_{Q,k}(u,z)\}$ 定义了层
$\tilde\pi_{E*}\cO_{\cC_E}(*S+k(\sigma'+\sigma''))$ 的一个局部全纯截面。这就证明了引理 4.16。\hfill$\square$

\begin{lemma}[4.17]
 设 $f\in\pi_*\cO_{\cC}(*S)$，并设 $\sum_{k=0}^{\infty}f_kq^k$ 为引理 4.16 中所定义的展开。则我们有
 \[
  \sum_{j=1}^{N}\langle\tilde\Psi|\sum_{k}\rho_j(X\otimes f_k)q^k=0.
 \]
 换言之，$\langle\tilde\Psi|$ 满足形式规范条件。
\end{lemma}

% [Tutor Note: 引理 4.17 的式子原文没有写出作用对象 $|\Phi\rangle$，是一个算子
%  恒等式：成立的是对所有 $|\Phi\rangle\in\cH^\dagger_\lambda(E)$ 有
%  $\sum_j\langle\tilde\Psi|\sum_k\rho_j(X\otimes f_k)q^k|\Phi\rangle=0$。
%  下面的证明正是补上 $|\Phi\rangle$ 后逐项验证，抄写时不要把它当成漏字删掉。]

\noindent\textit{证明.}\quad 按定义，对任意 $|\Phi\rangle\in\cH^\dagger_\lambda(E)$ 有
\begin{align*}
 \sum_{j=1}^{N}\langle\tilde\Psi|&\sum_{k=0}^{\infty}\rho_j(X\otimes f_k)q^k|\Phi\rangle\\
 &=\sum_{k=0}^{\infty}\sum_{d=0}^{\infty}\sum_{i=1}^{m_d}q^{k+d}\sum_{j=1}^{N}
   \langle\Psi|\rho_j(X\otimes f_k)|v_i(d)\otimes v^i(d)\otimes\Phi\rangle\\
 &=-\sum_{k=0}^{\infty}\sum_{d=0}^{\infty}\sum_{i=1}^{m_d}q^{k+d}
   \langle\Psi|\rho_{\sigma'}(X\otimes g_k)+\rho_{\sigma''}(X\otimes h_k)|v_i(d)\otimes v^i(d)\otimes\Phi\rangle.
\end{align*}
由 (4.27) 与 (4.28) 有
\begin{align*}
 \rho_{\sigma'}(X\otimes g_k)&=\sum_{m=0}^{\infty}f_{m,k}(t)\rho_{\sigma'}(X(m-k)),\\
 \rho_{\sigma''}(X\otimes h_k)&=\sum_{n=0}^{\infty}f_{k,n}(t)\rho_{\sigma''}(X(n-k)).
\end{align*}
由引理 4.14，我们有
\[
 (X(m-k)v_i(d)|v^j(d-m+k))+(v_i(d)|X(k-m)v^j(d-m+k))=0,
\]
于是
\begin{align*}
 &\sum_{i=1}^{m_d}\rho_{\sigma'}(X(m-k))|v_i(d)\otimes v^i(d)\otimes\Phi\rangle\\
 &\qquad+\sum_{j=1}^{m_{d-m+k}}\rho_{\sigma''}(X(-(m-k)))|v_j(d-m+k)\otimes v^j(d-m+k)\otimes\Phi\rangle=0.
\end{align*}
这就证明了引理 4.17。\hfill$\square$

% [Tutor Note: 证明第三行的负号与「两个点的分裂」是缝合的核心一步：把
%  $\langle\tilde\Psi|$ 按 (4.26) 展开成 $\sum_d q^d\langle\Psi_d|$、再按 (4.25)
%  收缩成基向量 $v_i(d)\otimes v^i(d)$ 之后，双点 $P$ 处的插入算子 $\rho_j$ 必须
%  拆成正规化后 $P'$（对应 $g_k$）与 $P''$（对应 $h_k$）两个贡献，而 $\rho_j$ 与
%  $\rho_{\sigma'}+\rho_{\sigma''}$ 相差一个负号 —— 这就是第三行前面那个 $-$ 的来源。
%  几何上即「把 $P$ 处的插入推到 $P'$ 与 $P''$ 处」。]

% [Tutor Note: (4.27)(4.28) 中的参数记作 $u$，到这里换成了 $t$；应为同一个局部
%  参数，属排版不一致，照原文抄。又 $X(m-k)$ 应读作仿射李代数元 $X\otimes t^{m-k}$
%  （即 $\rho_j$ 只吃掉 $f$ 的展开系数，幂次全部落到 $X$ 上）。]

% [Tutor Note: 末式两项相加为零正是引理 4.14 的不变性逐项作用的结果：第一项的
%  第一个因子取 $X(m-k)$、第二个因子取 $v^j(d-m+k)$，与第二项括号里的
%  $X(-(m-k))=X(k-m)$ 恰好构成 $(X(n)a|b)+(a|X(-n)b)=0$（$n=m-k$）。指标从
%  $v^i(d)$ 换成 $v^j(d-m+k)$ 是因为配对只在权相加为零的分量上非零。]

设 $\hat\cO_{\cB/E}$ 为沿 $E$ 的完备化，即
\[
 \hat\cO_{\cB/E}\simeq\varprojlim_{n}\cO_{\cB}/I_E^n,
\]
其中 $I_E$ 是 $E$ 的理想层。于是存在一个 $\cO_E$-代数同构
\[
 \hat\cO_{\cB/E}\simeq\cO_E[[q]].
\]
以下我们把 $\hat\cO_{\cB/E}$ 与 $\cO_E[[q]]$ 等同起来。

% [Tutor Note: 这里「沿 $E$ 的完备化同构于 $\cO_E[[q]]$」不是形式上的套话：在双点
%  附近 $\cB$ 的坐标是 $(u,q)$、而 $E$ 由 $q=0$ 定义，故 $I_E=(q)$，
%  $\cO_\cB/I_E^n$ 就是 $q$ 的 $n$ 次截断，取逆极限即得形式幂级数环。正是这个
%  同构让「形式规范条件」可以用 $q$ 的逐阶系数来表述（引理 4.16、4.17）。]

% [原文 p.86]

\begin{lemma}[4.18]
 设 $\hat\cV^\dagger_{\tilde\lambda/E}$ 为由形式规范条件所定义的 $\cO_E[[q]]$-模层：
 \[
  \hat\cV^\dagger_{\tilde\lambda/E}
  =\Big\{|\Phi\rangle\in\cH^\dagger_{\tilde\lambda}(E)[[q]]\ \Big|\
   \sum_{j=1}^{N}\langle\Phi|\sum_{k}\rho_j(X\otimes f_k)q^k=0
   \ \text{ 对一切 } f\in\pi_*\cO_{\cC}(*S)\Big\}.
 \]
 则存在一个 $\cO_E[[q]]$-模同构
 \begin{equation}
  \cV^\dagger_{\tilde\lambda}(\mathfrak F)\otimes_{\cO_\cB}\hat\cO_{\cB/E}
  \simeq\hat\cV^\dagger_{\tilde\lambda/E}. \tag{4.29}
 \end{equation}
\end{lemma}

% [Tutor Note: 定义中写的是 $\langle\Phi|$ 而不是引理 4.17 的 $\langle\tilde\Psi|$：
%  这里的元素是 $\cH^\dagger_{\tilde\lambda}(E)[[q]]$ 中的形式幂级数
%  $\sum_k\Phi_kq^k$，$\langle\Phi|$ 就是它的对偶向量；而 $\langle\tilde\Psi|$ 是
%  沿双点缝合出来、依赖 $q$ 的截面。两者不要混为一谈。]

\noindent\textit{证明.}\quad 由于 $\hat\cO_{\cB/E}$ 是忠实平坦的 $\cO_\cB$-模，我们有同构
\begin{align*}
 \cV^\dagger_{\tilde\lambda}(\mathfrak F)\otimes_{\cO_\cB}\hat\cO_{\cB/E}
 &=\underline{\Hom}_{\cO_\cB}(\cV_{\tilde\lambda}(\mathfrak F),\cO_\cB)\otimes_{\cO_\cB}\hat\cO_{\cB/E}\\
 &\simeq\underline{\Hom}_{\hat\cO_{\cB/E}}\big(\cV_{\tilde\lambda}(\mathfrak F)\otimes_{\cO_\cB}\hat\cO_{\cB/E},\hat\cO_{\cB/E}\big).
\end{align*}
注意我们有 $\cO_E$-模同构
\[
 \cH_{\tilde\lambda}(E)[[q]]\simeq\cH_{\tilde\lambda}\otimes_{\cO_\cB}\hat\cO_{\cB/E}.
\]
由引理 4.3 以及 $\hat\cO_{\cB/E}$ 的忠实平坦性，我们有
\[
 \cV_{\tilde\lambda}(\mathfrak F)\otimes_{\cO_\cB}\hat\cO_{\cB/E}
 \simeq\{\cH_{\tilde\lambda}\otimes_{\cO_\cB}\hat\cO_{\cB/E}\}
 \big/\{(\hat{\mathfrak g}(\mathfrak F)\cH_{\tilde\lambda})\otimes_{\cO_\cB}\hat\cO_{\cB/E}\}.
\]
另一方面我们有
\begin{align*}
 (\hat{\mathfrak g}(\mathfrak F)\cH_{\tilde\lambda})\otimes_{\cO_\cB}\hat\cO_{\cB/E}
 &=\hat{\mathfrak g}(\mathfrak F)\big(\cH_{\tilde\lambda}\otimes_{\cO_\cB}\hat\cO_{\cB/E}\big)\\
 &\simeq\hat{\mathfrak g}(\mathfrak F)\cH_{\tilde\lambda,E}[[q]],
\end{align*}
其中 $X\otimes f\in\hat{\mathfrak g}(\mathfrak F)$ 在 $\cH_{\tilde\lambda,E}[[q]]$ 上的作用由
\[
 \sum_{k=0}^{\infty}\sum_{j=1}^{N}\rho_j(X\otimes f_k)q^k
\]
给出，这里 $\sum f_kq^k$ 是引理 4.16 中所定义的与 $f$ 对应的元素。于是我们有
\[
 \cV_{\tilde\lambda}(\mathfrak F)\otimes_{\cO_\cB}\hat\cO_{\cB/E}
 \simeq\cH_{\tilde\lambda}(E)[[q]]\big/\hat{\mathfrak g}(\mathfrak F)\cH_{\tilde\lambda}(E)[[q]].
\]
这就证明了引理 4.18。\hfill$\square$

% [Tutor Note: 原书用带下划线的 $\underline{\Hom}$ 表示层 Hom（sheaf Hom），
%  以区别于整体截面的 $\Hom$；译文照排。第一行的第一个等号是定义
%  $\cV^\dagger=\underline{\Hom}_{\cO_\cB}(\cV,\cO_\cB)$，故 $\dagger$ 只在左端出现、
%  在 $\Hom$ 的变元里必须去掉，这不是漏排。]

% [Tutor Note: 「忠实平坦」（faithfully flat）在这里的作用是具体的：$\otimes_{\cO_\cB}
%  \hat\cO_{\cB/E}$ 正合，才能把引理 4.3 给出的 $\cV_{\tilde\lambda}$ 的商描述
%  原封不动地搬到完备化上（第二步），并把 $\hat{\mathfrak g}(\mathfrak F)$ 的作用
%  与张量积交换（第三步）。没有忠实平坦性，这两步都只是「形式上的」。]

% [Tutor Note: 末式左端仍写作 $\cV_{\tilde\lambda}$（没有 $\dagger$），而 (4.29) 的
%  左端是 $\cV^\dagger_{\tilde\lambda}$：原书的证明实际算的是 $\cV$ 的完备化
%  $\simeq$ 余不变量。要得到 (4.29) 还须再对偶一次 —— 商空间
%  $\cH_{\tilde\lambda}(E)[[q]]/\hat{\mathfrak g}(\mathfrak F)\cH_{\tilde\lambda}(E)[[q]]$
%  的对偶恰是 $\cH^\dagger_{\tilde\lambda}(E)[[q]]$ 中被该作用零化的子空间，即
%  定义中的 $\hat\cV^\dagger_{\tilde\lambda/E}$。原文把这一步整个省略了，读者需自己补上。]

\noindent\textbf{4.3.2 局部自由性.}\quad

\begin{theorem}[4.19]
 层 $\cV_{\tilde\lambda}(\mathfrak F)$ 是局部自由的 $\cO_\cB$-模。因此
 $\cV^\dagger_{\tilde\lambda}(\mathfrak F)$ 也是局部自由的 $\cO_\cB$-模。此外，对任意全纯映射
 $Y\to\cB$，设 $\mathfrak F_Y$ 为族 $\mathfrak F$ 的拉回。则我们有典范同构
 \[
  \cV^\dagger_{\tilde\lambda}(\mathfrak F)\otimes_{\cO_\cB}\cO_Y
  \simeq\cV^\dagger_{\tilde\lambda}(\mathfrak F_Y),
 \]
\end{theorem}

\noindent\textit{证明.}\quad 由推论 4.8，若我们的族 $\mathfrak F^{(0)}$ 不含奇异曲线（即 $D=\emptyset$），则定理成立。因此我们可以假设通用族 $\mathfrak F^{(0)}$ 含有奇异曲线。我们将对族 $\cC\to\cB$ 的双点轨迹的像的不可约分支的个数 $k$ 作归纳来证明本定理。

% [原文 p.87]

由于我们已假设 $k=1$，$\tilde{\mathfrak F}^{(0)}_E$ 是 $(N+2)$-点光滑曲线的通用族。选取足够小的 $\cB$，使得 $\bigoplus_{\mu\in P_\ell}\cV^\dagger_{\mu,\mu^\dagger,\tilde\lambda}(\tilde{\mathfrak F}_E)$ 是 $\cO_E$-自由的。设
\[
 \{\langle\Psi_1|,\dots,\langle\Psi_n|\}
\]
为 $\bigoplus_{\mu\in P_\ell}\cV^\dagger_{\mu,\mu^\dagger,\tilde\lambda}(\tilde{\mathfrak F}_E)$ 的一组 $\cO_E$-自由基。设
\[
 \{\langle\tilde\Psi_1|,\dots,\langle\tilde\Psi_n|\}
\]
为按 (4.26) 由 $\{\langle\Psi_1|,\dots,\langle\Psi_n|\}$ 构造出的 $\cV^\dagger_{\tilde\lambda}(\mathfrak F)$ 中的元素。对应 $\langle\Psi_i|\mapsto\langle\tilde\Psi_i|$（$i=1,\dots,n$）定义了 $\cO_E$-模同态
\[
 \iota:\bigoplus_{\mu\in P_\ell}\cV^\dagger_{\mu,\mu^\dagger,\tilde\lambda}(\tilde{\mathfrak F}_E)
 \to\hat\cV^\dagger_{\tilde\lambda/E}.
\]

首先我们证明 $\langle\tilde\Psi_1|,\dots,\langle\tilde\Psi_n|$ 是 $\cO_E[[q]]$-线性无关的。设有一关系
\[
 \sum_{i=1}^{n}a_i(q)\langle\tilde\Psi_i|=0,\qquad a_i(q)\in\cO_E[[q]].
\]
若对所有的 $i$ 都有 $a_i(0)=0$，则用 $q$ 的幂去除该关系，可设其中一个 $a_i(q)$（例如 $a_k(q)$）满足 $a_k(0)\ne0$。将上式令 $q=0$，得
\[
 \sum_{i=1}^{n}a_i(0)\langle\Psi_i|=0.
\]
由于 $\{\langle\Psi_1|,\dots,\langle\Psi_n|\}$ 是 $\cO_E$-线性无关的，故对所有的 $i$ 有 $a_i(0)=0$，矛盾。因此 $\langle\tilde\Psi_1|,\dots,\langle\tilde\Psi_n|$ 生成 $\hat\cV^\dagger_{\tilde\lambda/E}$ 的一个 $\cO_E[[q]]$-自由子模。

% [Tutor Note: 这里的「用 $q$ 的幂去除」之所以合法，是因为 $\cO_E[[q]]$ 是整环、$q$ 是非零因子：
%  先把关系按 $q$ 的最低次幂 $q^{m}$ 提出，再整体除以 $q^m$，就得到一个新的关系，
%  其中至少一个系数在 $q=0$ 处不为零（即 $a_k(0)\ne0$）。这一步是「$q$-进赋值」的
%  标准论证，漏掉它就看不出为什么可以假设 $a_k(0)\ne0$。]

现取一点 $x\in E$ 与 $s\in\cB\setminus E$。则由上述论证、推论 4.10 与 (4.29) 可得
\begin{align*}
 \dim_\C\cV_{\tilde\lambda}(\mathfrak F)\otimes\C_s
 &=\dim_\C\cV^\dagger_{\tilde\lambda}(\mathfrak F)\otimes\C_s\\
 &\ge n=\sum_{\mu\in P_\ell}\rank\cV^\dagger_{\mu,\mu^\dagger,\tilde\lambda}(\tilde{\mathfrak F}_E).
\end{align*}
由于 $\tilde{\mathfrak F}_E$ 在 $x$ 处是光滑曲线族，$\cV_{\mu,\mu^\dagger,\tilde\lambda}(\tilde{\mathfrak F}_E)$ 是局部自由的。由引理 4.3 与推论 4.9 我们有
\[
 \dim_\C\cV_{\tilde\lambda}(\mathfrak F)\otimes\C_x
 =\sum_{\mu\in P_\ell}\dim_\C\cV_{\mu,\mu^\dagger,\tilde\lambda}(\tilde{\mathfrak F}_E)\otimes\C_x.
\]
于是我们有
\[
 \dim_\C\cV_{\tilde\lambda}(\mathfrak F)\otimes\C_s
 \ge\dim_\C\cV_{\tilde\lambda}(\mathfrak F)\otimes\C_x.
\]
另一方面，由于 $\cV_{\tilde\lambda}(\mathfrak F)$ 在 $\cB\setminus E$ 上是凝聚且局部自由的，我们有不等式
\[
 \dim_\C\cV_{\tilde\lambda}(\mathfrak F)\otimes\C_s
 \le\dim_\C\cV_{\tilde\lambda}(\mathfrak F)\otimes\C_x.
\]
于是我们有等式
\[
 \dim_\C\cV_{\tilde\lambda}(\mathfrak F)\otimes\C_s
 =\dim_\C\cV_{\tilde\lambda}(\mathfrak F)\otimes\C_x.
\]
因此 $\cV_{\tilde\lambda}(\mathfrak F)$ 是局部自由的。\hfill$\square$

% [原文 p.88]

下面考虑 $k\ge2$ 的情形，此时判别轨迹有不可约分支分解 $D=\bigcup_{i=1}^{k}D_i$。设 $D_{sing}$ 为判别轨迹 $D$ 的奇点部分。上述论证表明 $\cV_{\tilde\lambda}(\mathfrak F)$ 在 $\cB\setminus D_{sing}$ 上是局部自由的。

\begin{lemma}[4.20]
 若对所有的 $i$，$\cV_{\tilde\lambda}(\mathfrak F_{D_i})$ 都是局部自由的 $\cO_{D_i}$-模，则 $\cV_{\tilde\lambda}(\mathfrak F)$ 是局部自由的 $\cO_\cB$-模。
\end{lemma}

\noindent\textit{证明.}\quad 令
\[
 \rank\cV_{\tilde\lambda}(\mathfrak F)|_{\cB\setminus D_{sing}}=m.
\]
则
\[
 \rank\cV_{\tilde\lambda}(\mathfrak F_{D_i})=m.
\]
因此对每个点 $x\in\cB$ 我们有
\[
 \dim_\C\cV_{\tilde\lambda}(\mathfrak F_x)=\dim_\C\cV_{\tilde\lambda}(\mathfrak F)\otimes\C_x=m.
\]
这说明，在每一点 $x\in\cB$ 处，$\cV_{\tilde\lambda}(\mathfrak F)$ 作为 $\cO_\cB$-模可由 $m$ 个截面在 $x$ 的邻域内生成，这是因为 $\cV_{\tilde\lambda}(\mathfrak F)$ 是凝聚的 $\cO_\cB$-模。另一方面，若 $\cV_{\tilde\lambda}(\mathfrak F)$ 在某点 $y\in D_{sing}$ 处不是局部自由的，则作为 $\cO_\cB$-模，$\cV_{\tilde\lambda}(\mathfrak F)$ 在点 $y$ 的邻域内需要多于 $m$ 个截面才能生成，这仍是因为 $\cV_{\tilde\lambda}(\mathfrak F)$ 是凝聚的 $\cO_\cB$-模。因此 $\cV_{\tilde\lambda}(\mathfrak F)$ 是局部自由的。这就证明了引理 4.20。

于是证明归结为证明层 $\cV_{\tilde\lambda}(\mathfrak F_{D_i})$ 的局部自由性。为此令 $E=D_i$，并考虑如下映射：
\[
\begin{array}{ccc}
 j_E:\cV_{\tilde\lambda}(\mathfrak F_E) & \longrightarrow & \bigoplus\limits_{\mu\in P_\ell}\cV_{\mu,\mu',\tilde\lambda}(\tilde{\mathfrak F}_E)\\[3pt]
 \mbox{}[|\phi\rangle] & \longmapsto & ([\,|0_{\mu,\mu^\dagger}\otimes\phi\rangle\,])
\end{array}
\]
其中 $[|\phi\rangle]$（相应地，$[|0_{\mu,\mu^\dagger}\otimes\phi\rangle]$）是 $|\phi\rangle\in\cH_{\tilde\lambda}(E)$（相应地，$|0_{\mu,\mu^\dagger}\otimes\phi\rangle\in\cH_{\mu,\mu',\tilde\lambda}(E)$）在 $\cV_{\tilde\lambda}(\mathfrak F_E)$（相应地，$\cV_{\mu,\mu',\tilde\lambda}(\tilde{\mathfrak F}_E)$）中的像。\hfill$\square$

\begin{lemma}[4.21]
 映射 $j_E$ 是定义良好的，并且是满射。此外，若对每个 $\mu\in P_\ell$，$\cV_{\mu,\mu',\tilde\lambda}(\tilde{\mathfrak F}_E)$ 都是局部自由的 $\cO_E$-模，则 $j_E$ 是同构，从而是局部自由的。
\end{lemma}

\noindent\textit{证明.}\quad 设元素 $|\phi\rangle$ 与 $|\phi'\rangle$ 满足关系
\[
 |\phi'\rangle=|\phi\rangle+\sum_\nu X_\nu\otimes f_\nu|\psi_\nu\rangle,
\]
其中
\[
 X_\nu\otimes f_\nu\in\hat{\mathfrak g}(\mathfrak F_E),\qquad |\psi_\nu\rangle\in\cH_{\tilde\lambda}(E).
\]
设 $\tilde f_\nu$ 为 $f_\nu$ 到 $\tilde{\mathfrak F}_E$ 的拉回。由于 $\tilde f_\nu$ 在 $\sigma'(E)$ 与 $\sigma''(E)$ 处全纯，我们有
\[
 \rho_{\sigma'(E)}(X_\nu\otimes\tilde f_\nu)|0_{\mu,\mu^\dagger}\rangle
 +\rho_{\sigma''(E)}(X_\nu\otimes\tilde f_\nu)|0_{\mu,\mu^\dagger}\rangle=0.
\]
于是我们有
\[
 |0_{\mu,\mu^\dagger}\otimes\phi'\rangle=|0_{\mu,\mu^\dagger}\otimes\phi\rangle
 +\sum_\nu X_\nu\otimes\tilde f_\nu|0_{\mu,\mu^\dagger}\otimes\psi_\nu\rangle.
\]
因此我们有
\[
 [\,|0_{\mu,\mu^\dagger}\otimes\phi\rangle\,]=[\,|0_{\mu,\mu^\dagger}\otimes\phi'\rangle\,].
\]
从而映射 $j_E$ 是定义良好的。

% [原文 p.89]

其次我们考虑正合列
\[
 \cV_{\tilde\lambda}(\tilde{\mathfrak F}_E)\xrightarrow{\ j_E\ }
 \bigoplus_{\mu\in P_\ell}\cV_{\mu,\mu^\dagger,\tilde\lambda}(\tilde{\mathfrak F}_E)
 \to\operatorname{Coker}j_E\to0.
\]
由推论 3.20 有同构
\[
 \cV_{\tilde\lambda}(\tilde{\mathfrak F}_E)\otimes\C_x
 \simeq\bigoplus_{\mu\in P_\ell}\cV_{\mu,\mu^\dagger,\tilde\lambda}(\tilde{\mathfrak F}_E)\otimes\C_x.
\]
故 $\operatorname{Coker}j_E\otimes\C_x=0$。于是由中山引理（Nakayama's lemma）有 $\operatorname{Coker}j_E=0$。从而 $j_E$ 是满射。最后考虑正合列
\[
 0\to\operatorname{Ker}j_E\to\cV_{\tilde\lambda}(\tilde{\mathfrak F}_E)
 \to\bigoplus_{\mu\in P_\ell}\cV_{\mu,\mu^\dagger,\tilde\lambda}(\tilde{\mathfrak F}_E)\to0.
\]
假设 $\bigoplus_{\mu\in P_\ell}\cV_{\mu,\mu^\dagger,\tilde\lambda}(\tilde{\mathfrak F}_E)$ 是局部自由的。则我们有
\[
 \underline{\operatorname{Tor}}^1_{\cO_E}
 \Big(\bigoplus_{\mu\in P_\ell}\cV_{\mu,\mu^\dagger,\tilde\lambda}(\tilde{\mathfrak F}_E),\C_x\Big)=0.
\]
于是我们得到正合列
\[
 0\to\operatorname{Ker}j_E\otimes\C_x\to\cV_{\tilde\lambda}(\tilde{\mathfrak F}_E)\otimes\C_x
 \to\bigoplus_{\mu\in P_\ell}\cV_{\mu,\mu^\dagger,\tilde\lambda}(\tilde{\mathfrak F}_x)\otimes\C_x\to0.
\]
因此 $\operatorname{Ker}j_E\otimes\C_x=0$。于是再由中山引理可得 $\operatorname{Ker}j_E=0$。故 $j_E$ 是同构。这就证明了引理 4.21。\hfill$\square$

最后我们对 $(g,N)$ 作归纳来证明定理 4.19。当 $g=0$、$N=3$ 时，$\cB$ 是一个点。故此时定理成立。设定理对 $g=0$ 与 $N\le M$ 成立。令 $N=M+1$。则对亏格 0 的 $N$-点稳定曲线的局部通用族 $\mathfrak F^{(0)}$，有 $D\ne\emptyset$，并且 $\tilde{\mathfrak F}^{(0)}_{D_i}$ 是若干个 $(N'+1)$-点与 $(N''+1)$-点亏格 0 稳定曲线的局部通用族的无交并，其中 $N'+N''=N$。由于 $N'+1\le M$、$N''+1\le M$，由归纳假设，$\cV_{\mu,\mu^\dagger,\tilde\lambda}(\tilde{\mathfrak F}_E)$ 是局部自由的。于是由上述论证，定理成立。这样我们就能证明定理对任意正整数 $N$ 的 $(0,N)$ 成立。

最后设定理对任意满足 $g\le h$、$N\le M$ 的 $(g,N)$ 成立。则由与上面类似的论证，定理对 $(g,M+1)$ 成立。设 $\mathfrak F^{(0)}$ 是亏格 $h+1$ 的 $N$-点稳定曲线的局部通用族，其中 $N\le M$。若该族的判别轨迹 $D$ 为空，即族中每条曲线都非奇异，则由推论 4.8 定理成立。若 $D\ne\emptyset$，则 $\tilde{\mathfrak F}^{(0)}_{D_i}$ 要么是亏格 $h$ 的 $(N+2)$-点稳定曲线的局部通用族，要么是亏格 $g'$ 的 $(N'+1)$-点稳定曲线与亏格 $g''$ 的 $(N''+1)$-点稳定曲线的局部通用族的无交并，其中 $N'+N''=N$、$g'+g''=h$。若 $g'<h$ 且 $g''<h$，则由归纳假设定理成立。若 $g'=0$，则 $N'\ge2$，因为该族是亏格 0 的 $(N'+1)$-点稳定曲线。于是 $g''=h$ 且 $N''+1<N$。因此再由归纳假设，定理成立。这就完全证明了定理 4.19。

上述引理 4.21 还有一个重要的推论。为陈述它，我们稍微推广一下情形。令
\[
 E=\bigcap_{i=1}^{N}D_i.
\]

% [原文 p.90]

设 $\mathfrak F^{(0)}_E$ 为我们族的拉回，并设 $\tilde{\mathfrak F}^{(0)}_E$ 为沿双点作正规化所得的族。设 $\sigma'_i,\sigma''_i$ 为相应于轨迹 $D_i$（$i=1,\dots,k$）的双点处的全纯截面。我们还引入 $\cB$ 的局部坐标 $(u_1,\dots,u_{M-k},q_1,\dots,q_k)$，使得 $D_i$ 由 $q_i=0$ 定义。此外，在 $D_i$ 上方双点轨迹的邻域内，我们选取 $\cC$ 的局部坐标 $(u_1,\dots,u_{M-k},q_1,\dots,q_{i-1},z_i,w_i,q_{i+1},\dots,q_k)$，使得态射 $\pi$ 由
\[
 (u_1,\dots,u_{M-k},q_1,\dots,q_{i-1},z_i,w_i,q_{i+1},\dots,q_k)
 \mapsto(u_1,\dots,u_{M-k},q_1,\dots,q_{i-1},z_i\cdot w_i,q_{i+1},\dots,q_k)
\]
给出。然后，在 $\sigma'(E)$（相应地，$\sigma_i''(E)$）的邻域内，我们选取 $\tilde{\mathfrak F}^{(0)}_E$ 的局部坐标 $(u_1,\dots,u_k,z_i)$（相应地，$(u_1,\dots,u_k,w_i)$），并令
\begin{align*}
 \tilde{\mathfrak F}_E:={}&(\pi_E:\cC_E\to E;s_1|_E,\dots,s_N|_E,
 \sigma'_1,\dots,\sigma'_k,\sigma''_1,\dots,\sigma''_k;\\
 &\eta_1|_{\cC_E},\dots,\eta_N|_{\cC_E},z_1,\dots,z_k,w_1,\dots,w_k).
\end{align*}
有了这些记号，引理 4.21 蕴涵下述定理。

\begin{theorem}[4.22]
 （因子化）存在 $\cO_E$-模同构
 \[
  \bigoplus_{\vec\mu\in P_\ell^k}\cV^\dagger_{\vec\mu,\vec\mu^\dagger,\tilde\lambda}(\tilde{\mathfrak F}_E)
  \simeq\cV^\dagger_{\tilde\lambda}(\mathfrak F_E),
 \]
 其中，对 $\vec\mu=(\mu_1,\dots,\mu_k)$，$\vec\mu^\dagger$ 由 $\vec\mu^\dagger=(\mu_1^\dagger,\dots,\mu_k^\dagger)$ 定义。
\end{theorem}

% [Tutor Note: 这个「因子化」定理就是本节标题所承诺的结论，值得记住它的几何含义：
%  $E=\bigcap_i D_i$ 是判别轨迹中最深的层，在 $E$ 上方曲线同时有 $k$ 个双点；把每个
%  双点沿 $\sigma'_i,\sigma''_i$ 分开就得到 $\tilde{\mathfrak F}_E$。定理说，$E$ 上的
%  共形块层可以按「在每个双点处插入一组完备中间态」来分解：直和跑遍 $\vec\mu\in
%  P_\ell^k$，每个分量 $\cV^\dagger_{\vec\mu,\vec\mu^\dagger,\tilde\lambda}$ 是带
%  $2k$ 个额外插入点的共形块层。$k=1$ 时正是引理 4.21 的 $j_E$ 为同构的结论；
%  一般的 $k$ 由对 $k$ 归纳（逐个分开双点）得到。]
